% !TEX root = ../main.tex

The Expedia Personalised Hotel Searches dataset originates from the ICDM 2013 competition of the same name and is reused for the in-class Kaggle competition of the VU MSc Data Mining Techniques course. Each row of the training and test files corresponds to one (search, property) pair, and the task is to rank the candidate hotels within each search query so that booked and clicked properties appear at the top. The course evaluates submissions with NDCG@5 under the IDCG = 0 convention, which scores searches with no positive labels as 0 rather than leaving them undefined. This convention has a direct impact on the headline number, because roughly half of the holdout searches contain no booking or click in the top five even at the model's best operating point.

The task fits the listwise learning-to-rank framing. LambdaMART, introduced by Burges~\cite{burges2010ranknet} as an extension of MART with a listwise gradient that approximates NDCG, is the natural model family and is also what the top-ranked entries of the 2013 competition used. Liu et al.~\cite{liu2013expedia} combined an ensemble of about thirty base learners, with LambdaMART carrying the largest weight, to reach the 1st-place private leaderboard score of NDCG@38 = 0.535. Wang and Kalousis~\cite{wang2014expedia} placed second with roughly 300 engineered features dominated by within-query relativisations of price and quality. More directly comparable to our setting, because the metric is NDCG@5 on the in-class leaderboard rather than NDCG@38 on the ICDM private set, is Bellucci et al.~\cite{bellucci2021dmt}, whose 2021 VU MSc Data Mining Techniques course report reported a public NDCG@5 of 0.41726 using LightGBM LambdaMART with leave-one-out (LOO) target encoding on \texttt{prop\_id} and a per-property positional proxy from rows where Expedia had shuffled the page (\texttt{random\_bool} = 1).

Our pipeline draws on this prior work but is otherwise built from scratch as six numbered notebooks. Three contributions are worth flagging up front. First, target encoding follows the Bellucci recipe but uses 5-fold out-of-fold (OOF) encoding rather than LOO; OOF is computationally cheaper at 4.96M rows and avoids the Bellucci variance at low-impression \texttt{prop\_id} values. Second, hyperparameter tuning uses Optuna's Tree-structured Parzen Estimator (TPE)~\cite{bergstra2011tpe,akiba2019optuna} on both LightGBM and XGBoost, with the LightGBM ensemble blended with XGBoost at a val-chosen weight. Third, a deployment-time fairness-of-exposure rerank, following the Singh and Joachims~\cite{singh2018fairness} framework, is added on top of the trained models and tuned on val with a pre-specified accuracy-fairness contract before being measured on the locked holdout.

The remainder is structured as follows. Section~\ref{sec:data} covers exploration, cleaning, and feature engineering. Section~\ref{sec:models} describes algorithm choice, Optuna tuning, the position-bias correction experiment, and the seeded ensemble plus LightGBM-XGBoost blend. Section~\ref{sec:eval} reports holdout NDCG@5 by predictor, by search-page size, and by ranked vs random page, together with a worst-case characterisation and feature importance. Section~\ref{sec:bias} presents the bias audit and the inference-time rerank. Section~\ref{sec:final} summarises the final pipeline and reports the realised Kaggle public score. Section~\ref{sec:disc} closes with a discussion that includes a short reflection on NDCG@5 versus alternative ranking metrics, limitations, and reproducibility.
